

\subsection{The Standard Model and Dark Matter}

Explain the standard model briefly and introduce dark matter and connect the two as context for the reader.

\subsection{The Light Dark Matter Experiment (LDMX)}

Introduce the Light Dark Matter Experiment as a bridge from the last subchapter

Give some practical info and history

Explain the research goal and what is being investigated


\newpage

\subsection{Motivation}
Introduce the physics motivation behind LDMX.
Explain the importance of precise electron counting and pile-up mitigation and how that relates to the experiment's \textit{physics reach}.
Motivate why traditional clustering methods may struggle in multi-electron environments.
Introduce graph neural networks (GNNs) as a promising approach for irregular detector data.

(!) This has already been very nicely done by others. Take inspiration from earlier thesis reports and make one that takes the best things from all of them.

\subsection{Research Objectives}
Clearly define the problem: robust electromagnetic shower separation under pile-up conditions.
State the main research questions.
Define performance goals and comparison to baseline methods (e.g., CLUE, ParticleFlow).

\subsection{Thesis Structure}
Briefly describe how the thesis is organized.




\subsection{The LDMX Experiment}
Describe the purpose of LDMX.
Explain relevant detector components: ECal, HCal, tracking system.
Describe the simulation framework (ldmx-sw, Geant4). Explain dark bremsstrahlung.

(!) This has already been very nicely done by others. Take inspiration from earlier thesis reports and make one that takes the best things from all of them.

*Figure LDMX detector overview*

*Figure dark bremsstrahlung particle interaction*

\subsection{Electromagnetic Showers}
Explain shower formation in calorimeters.
Discuss spatial and energy deposition characteristics.
Describe challenges arising from overlapping showers.

*Figure ECal cells and trigger cells*

\subsection{Rule-Based Reconstruction in High Energy Physics}
Introduce clustering and ParticleFlow.
Describe existing LDMX reconstruction approaches (e.g., CLUE).
Discuss online vs offline reconstruction context in order to give the reader an idea of online/offline and where this reconstruction method aims to relate.

*Figure that explains CLUE and PF*


\subsection{Graph Neural Networks}
Introduce graph representations and message passing.
Explain why calorimeter data can naturally be represented as graphs. Explain how data from other subdetectors can be included in the same graph representation and refer to MLPF by CMS.
Explain attention mechanism in GNN.
Describe relevant architectures (e.g., MLPF, GravNet-style, ParticleNet-inspired).
Discuss advantages compared to CNNs/RNNs for irregular geometries.
Discuss the prospects of a transformer based architecture and explain the analogy between transformers and GNN (especially when GNN has attention). Perhaps refer to entirely transformer based solutions presented by CMS in MLPF. 

*Figure flowchart of where the different alternatives of reconstruction (Rule based CLUE, PF and ML) fit in the reconstruction pipeline of LDMX*

*Figure illustrating the idea of GNN in HEP reconstruction and subdetector data (maybe omitt if there is a proper one later anyways)*
